\documentclass[journal]{IEEEtran}

\usepackage[T1]{fontenc}
\usepackage{amsmath,amsfonts}
\usepackage{newtxtext}
\usepackage{newtxmath}

\usepackage{algorithmic}

\usepackage{cite}
\usepackage{graphicx}
\usepackage{booktabs}

\usepackage[hidelinks]{hyperref}

\begin{document}

\title{title}

\author{Andy Babcock \\
    \textit{Steve Sanghi College of Engineering, Northern Arizona University} \\
    Flagstaff, AZ, USA \\
    aab726@nau.edu
}

\maketitle
What is the effect of noise on test-taking anxiety and performance?

Listening to white noise shows no statistically significant reduction in anxiety or improvement in test scores, tested with a wearable EEG measuring anxiety. If there is a reduction in anxiety, there may also be an improvement in test scores. Additionally, we will test familiar and unfamiliar music, with the expectation that familair music reduces anxiety and improves test scores, while unfamiliar music reduces anxiety at the cost of lower test scores. Alternatively, familiar music may not improve test scores, or unfamilir music may improve test scores similarly to familiar music.

\bibliographystyle{IEEEtran}
\bibliography{references}

\end{document}
